\section{Stage 3A: solver-relative hardness tracking}
\label{sec:hardness-tracking}

This section studies whether the verified HGB40 plus Elite-QD1000
$(P=T)$ mechanism tracks high current BQ hardness inside each fixed niche as
the solver changes.  It does not identify current hardness with future
reachability or downstream training value, and it does not compare BLADE to
HAC.  All maxima below are over finite realized candidate sets; no finiteness
of the ambient instance space is assumed.

\subsection{Tracking objects and information order}
Let $M_k$ be the BQ solver entering cycle $k$ and define, on the same scale as
the fitted HGB target,
\[
 f_k(x)=\log(1+L_{\rm BQ}(M_k,x)).
\]
Let $\widehat f_k$ be the median HGB40 prediction.  For every
$j\in\mathsf S(C_k)$ write
\[
 E_{k,j}=C_k\cap\mathcal N_j,
 \qquad
 e_{k,j}=\mathop{\arg\max}_{x\in E_{k,j}}^{\rm score/hash}
                 \widehat f_k(x),
\]
where ``score/hash'' means descending prediction followed by ascending
instance hash.  Thus $e_{k,j}$ is a single, deterministic member of the
per-niche floor $Q_k$, and
\[
 e_{k,j}\in Q_k\subseteq T_k=P_{k+1}.
\]
The prediction and selection use $M_k$.  Exact labels on $T_k$ and adaptation
to $M_{k+1}$ occur only after this selection.  Any inequality involving
$f_{k+1}(e_{k,j})$ is therefore an ex post tracking statement under a declared
drift premise, not a claim that cycle $k$ used future information.

The quantities kept distinct in this section are: current true hardness
$f_k$, current prediction $\widehat f_k$, within-niche selected quality,
solver-induced drift $f_{k+1}-f_k$, best hardness among all historically
generated candidates, future mutation reachability, and downstream BQ
training value.  No implication among the last three is assumed.

\subsection{One-step approximate within-niche maximization}
\mstatus{PROVED under a finite-pool error bound}
\istatus{PARTIAL} \nstatus{STANDARD}

\begin{proposition}[Per-niche approximate argmax]
\label{prop:tracking-argmax}
Fix $k$ and $j\in\mathsf S(C_k)$.  Suppose $\epsilon_{k,j}\geq0$ and
\[
 \lvert\widehat f_k(x)-f_k(x)\rvert\leq\epsilon_{k,j}
 \qquad\text{for every }x\in E_{k,j}.
\]
Then
\[
 f_k(e_{k,j})\geq
 \max_{x\in E_{k,j}}f_k(x)-2\epsilon_{k,j}.
\]
If $P_k\cap\mathcal N_j\neq\varnothing$, parent inclusion in the candidate
union also gives
\[
 f_k(e_{k,j})\geq
 \max_{x\in P_k\cap\mathcal N_j}f_k(x)-2\epsilon_{k,j}.
\]
\end{proposition}

\begin{proof}
Because $E_{k,j}$ is finite and nonempty, choose
$x^*\in\arg\max_{x\in E_{k,j}}f_k(x)$.  The selected candidate maximizes the
prediction even when the prediction maximum is tied, so
\begin{align*}
 f_k(e_{k,j})
 &\geq \widehat f_k(e_{k,j})-\epsilon_{k,j}\\
 &\geq \widehat f_k(x^*)-\epsilon_{k,j}\\
 &\geq f_k(x^*)-2\epsilon_{k,j}.
\end{align*}
The parent-relative result follows from
$P_k\cap\mathcal N_j\subseteq E_{k,j}$.
\end{proof}

The deterministic hash rule makes the selected maximizer single-valued but
does not change the proof.  An empty candidate niche has no elite and no
maximum, so the statement is restricted to $j\in\mathsf S(C_k)$.  The
parent-relative form additionally requires that the parent niche be nonempty.
Hash deduplication must preserve the phenotype, score, and niche identity used
by the statement on the realized pool.  The additive $2\epsilon$ guarantee is
on the $\log(1+\mathrm{LBQ})$ scale.  On the raw scale it implies only
\[
 L_{\rm BQ}(M_k,e_{k,j})
 \geq e^{-2\epsilon_{k,j}}
       \left(1+\max_{x\in E_{k,j}}L_{\rm BQ}(M_k,x)\right)-1,
\]
not an additive raw-LBQ bound of the same size.

This is the standard approximate-argmax inequality.  Elite-QD makes the
chosen approximate maximizer survive into $P_{k+1}$; that implementation fact
is useful, but it does not make the inequality itself novel.  S6-AE verifies
the score scale, finite candidate pool, score/hash rule, and $P=T$.  It does
not verify the uniform error premise on $E_{k,j}$, because exact current LBQ
is not obtained for the complete candidate union.

\subsection{One-step tracking across solver adaptation}
\mstatus{PROVED under local error and drift bounds}
\istatus{PARTIAL/OPEN} \nstatus{STANDARD}

\begin{proposition}[Drift-adjusted per-niche tracking]
\label{prop:tracking-drift}
Under Proposition~\ref{prop:tracking-argmax}, suppose additionally that
\[
 \lvert f_{k+1}(x)-f_k(x)\rvert\leq\delta_{k,j}
 \qquad\text{for every }x\in E_{k,j},
\]
where $\delta_{k,j}\geq0$.  Then
\[
 f_{k+1}(e_{k,j})\geq
 \max_{x\in E_{k,j}}f_{k+1}(x)
 -2\epsilon_{k,j}-2\delta_{k,j}.
\]
Consequently, whenever $P_k\cap\mathcal N_j$ is nonempty,
\[
 f_{k+1}(e_{k,j})\geq
 \max_{x\in P_k\cap\mathcal N_j}f_{k+1}(x)
 -2\epsilon_{k,j}-2\delta_{k,j}.
\]
For only the parent-relative conclusion it is sufficient that the drift bound
hold on $(P_k\cap\mathcal N_j)\cup\{e_{k,j}\}$.
\end{proposition}

\begin{proof}
The drift bound and Proposition~\ref{prop:tracking-argmax} give
\begin{align*}
 f_{k+1}(e_{k,j})
 &\geq f_k(e_{k,j})-\delta_{k,j}\\
 &\geq \max_{x\in E_{k,j}}f_k(x)
           -2\epsilon_{k,j}-\delta_{k,j}\\
 &\geq \max_{x\in E_{k,j}}f_{k+1}(x)
           -2\epsilon_{k,j}-2\delta_{k,j}.
\end{align*}
Restricting the comparison set to old parents proves the second claim.  The
same calculation uses only the smaller stated set for the parent-relative
version.
\end{proof}

The guarantee says that an elite selected under $M_k$ remains competitive,
after one update, with candidates that were already available in cycle $k$.
It does not compare against new cycle-$(k+1)$ offspring, show that
$e_{k,j}$ remains the next HGB elite, preserve its individual hash beyond
$P_{k+1}$, or control future reachability or training value.  The drift bound
is local to an explicitly named finite set; no global bound on the measurable
TSP space is needed.  Even this local bound is not currently certified by
S6-AE.

\subsection{Historical benchmark and multicycle regret}
For a finite horizon define the latent generated-candidate ledger
\[
 \mathcal U_k=\bigcup_{t=0}^{k}C_t.
\]
For $j\in\mathsf S(\mathcal U_k)$ define
\[
 V_{k,j}=\max_{x\in\mathcal U_k\cap\mathcal N_j}f_k(x),
 \qquad
 R_{k,j}=V_{k,j}-f_k(e_{k,j}).
\]
By deterministic niche persistence, once niche $j$ has appeared by cycle
$k$, it is occupied in $P_{k+1}$ and hence has an elite in every later cycle.
The benchmark $V_{k,j}$ is theoretical: S6-AE neither retains every member of
$\mathcal U_k$ as a candidate nor evaluates every historical candidate under
$M_k$.

Define the current unavailability penalty
\[
 A_{k,j}=\left[
 \max_{x\in\mathcal U_{k-1}\cap\mathcal N_j}f_k(x)
 -\max_{x\in C_k\cap\mathcal N_j}f_k(x)
 \right]_+,
\]
with $A_{k,j}=0$ if the historical set is empty.

\begin{proposition}[Exact source of historical tracking loss]
\label{prop:tracking-decomposition}
Under the prediction-error premise of
Proposition~\ref{prop:tracking-argmax},
\[
 R_{k,j}\leq2\epsilon_{k,j}+A_{k,j}.
\]
\end{proposition}

\begin{proof}
Write $B_{k,j}=\max_{x\in C_k\cap\mathcal N_j}f_k(x)$.  Since
$\mathcal U_k=\mathcal U_{k-1}\cup C_k$,
$V_{k,j}\leq B_{k,j}+A_{k,j}$.  Proposition
~\ref{prop:tracking-argmax} gives
$f_k(e_{k,j})\geq B_{k,j}-2\epsilon_{k,j}$, and subtraction proves the claim.
\end{proof}

This non-accumulating decomposition is informative only if $A_{k,j}$ can be
controlled.  Parent--child union makes the immediately previous elite
available, which yields the following recursive control under a much stronger
historical drift premise.

\mstatus{PROVED under full historical-set drift}
\istatus{OPEN} \nstatus{STANDARD}

\begin{proposition}[Historical per-niche regret recurrence]
\label{prop:tracking-recurrence}
Suppose niche $j$ was occupied by cycle $k-1$.  Assume:
\begin{enumerate}
\item the prediction-error premise holds on $C_k\cap\mathcal N_j$;
\item $e_{k-1,j}\in P_k\subseteq C_k$; and
\item for every $x\in\mathcal U_{k-1}\cap\mathcal N_j$,
\[
 \lvert f_k(x)-f_{k-1}(x)\rvert\leq\delta_{k-1,j}.
\]
\end{enumerate}
Then
\[
 R_{k,j}\leq R_{k-1,j}
       +2\epsilon_{k,j}+2\delta_{k-1,j}.
\]
If $\tau_j$ is the first cycle in which niche $j$ appears, then
\[
 R_{\tau_j,j}\leq2\epsilon_{\tau_j,j},
\]
and for every $k\geq\tau_j$,
\[
 R_{k,j}\leq
 2\sum_{s=\tau_j}^{k}\epsilon_{s,j}
 +2\sum_{s=\tau_j}^{k-1}\delta_{s,j}.
\]
\end{proposition}

\begin{proof}
Let $B_{k,j}=\max_{C_k\cap\mathcal N_j}f_k$.  The retained prior elite gives
$B_{k,j}\geq f_k(e_{k-1,j})$.  Uniform drift over the historical set gives
\begin{align*}
 \max_{x\in\mathcal U_{k-1}\cap\mathcal N_j}f_k(x)
 &\leq V_{k-1,j}+\delta_{k-1,j}\\
 &=f_{k-1}(e_{k-1,j})+R_{k-1,j}+\delta_{k-1,j}\\
 &\leq f_k(e_{k-1,j})+R_{k-1,j}+2\delta_{k-1,j}\\
 &\leq B_{k,j}+R_{k-1,j}+2\delta_{k-1,j}.
\end{align*}
New candidates in $\mathcal U_k\setminus\mathcal U_{k-1}$ lie in $C_k$, so
their maximum is at most $B_{k,j}$.  Hence
$V_{k,j}\leq B_{k,j}+R_{k-1,j}+2\delta_{k-1,j}$.  The approximate-argmax
bound $f_k(e_{k,j})\geq B_{k,j}-2\epsilon_{k,j}$ proves the recurrence.

At first discovery there is no older candidate in niche $j$, so
$\mathcal U_{\tau_j}\cap\mathcal N_j=C_{\tau_j}\cap\mathcal N_j$ and the
initial bound follows directly.  Iterating the recurrence gives the stated
finite-horizon sum.
\end{proof}

The displayed recurrence proposed for Stage 3A is therefore valid, but only
under drift control on the \emph{entire historical niche ledger}, including
candidates that the algorithm discarded and never relabels.  Bounded
per-cycle errors and drift yield a bound that grows linearly with the horizon;
they do not give bounded long-run regret.  A steady-state statement would
need, for example, summable error and drift, a uniform bound on
$A_{k,j}$, or an operational archive that repeatedly reintroduces the
historical benchmark candidates.  None is verified for S6-AE.

The accumulation cannot generally be removed.  In a single niche with
capacity one and zero solver drift, let cycle $t\geq1$ contain the retained
$x_{t-1}$ and a new $x_t$, with $f(x_t)=-2t\epsilon$.  Give both candidates
the tied prediction $-(2t-1)\epsilon$ by errors of magnitude $\epsilon$, and
choose hashes so that the deterministic tie rule selects $x_t$.  Starting
from $f(x_0)=0$, the historical regret after cycle $t$ is exactly
$2t\epsilon$.  Thus a constant error bound alone permits linear loss even
with parent--child union and no solver drift.

\subsection{What happens after all 1,000 niches are occupied}
Proposition~\ref{prop:niche-persistence} makes structural coverage constant
once every niche is occupied.  Propositions
~\ref{prop:tracking-argmax}--\ref{prop:tracking-recurrence} address a different
quantity: which representative is selected and how close its current
solver-relative hardness is to a declared comparator.  The per-niche floor
continues to replace representatives after coverage saturation.  Therefore
coverage can be saturated while within-niche hardness quality improves,
degrades after solver adaptation, or loses future stepping stones.  Global
HGB fill can retain additional candidates but supplies no worst-case
historical or reachability certificate.

\subsection{Adversarial boundary checks}
The following finite constructions are exact mathematical checks, not claims
about the realized S6-AE trajectory.

\paragraph{A. Perfect prediction with severe solver drift.}
Put $a,b$ in one niche and let $P_k=C_k=\{a,b\}$.  Set
$f_k(a)=1$, $f_k(b)=0$, and $\widehat f_k=f_k$, so $e_{k,j}=a$.  After
adaptation set $f_{k+1}(a)=0$ and $f_{k+1}(b)=H$ for arbitrary $H>0$.
The post-update gap to the old-parent maximum is $H$.  Perfect current
prediction gives no cross-checkpoint guarantee without a drift bound.

\paragraph{B. Small drift with inaccurate prediction.}
Keep the solver fixed, so $\delta=0$.  Set $f(a)=H$, $f(b)=0$ in one niche,
but $\widehat f(a)=0$ and $\widehat f(b)=1$.  The elite is $b$ and its current
quality gap is $H$, which can be arbitrarily large.  Small solver drift cannot
replace a candidate-set prediction or ranking guarantee.

\paragraph{C. Full coverage with deteriorating future reachability.}
The progressive bridge-loss construction in
Section~\ref{sec:progressive-bridge-loss} retains every occupied niche while
successively replacing every future-useful bridge by a higher-current-
hardness decoy.  It applies even with perfect current-hardness prediction and
the complete global-fill mechanism.  Current-hardness tracking therefore does
not imply future-reachability tracking.

\paragraph{D. Zero tracking regret with no downstream training benefit.}
Let every cycle have a unique current-hardness maximizer $a_k$ in its niche
and let selection be exact, so $R_{k,j}=0$.  Define the solver adaptation map
on the selected training set to leave the solver unchanged, or to move it to
a state with worse loss on a separate evaluation instance $z$.  Then tracking
regret is identically zero while downstream performance is unchanged or
worse.  A link from selected hardness to the optimization effect of a BQ
update is an additional assumption, not a consequence of elite tracking.

\paragraph{E. A reset method can win without elite preservation.}
Let the inherited population contain only $a$, let all scheduled mutations
from $a$ return $a$, and let the current hard target be $\{g\}$.  A reset
generator that returns $g$ with probability one hits in one equal-cost
opportunity; the inherited method never hits.  Historical elite preservation
does not dominate a reset or stationary proposal law with better direct
target mass.

\subsection{Applicability and novelty audit}
\begin{center}
\begin{tabular}{>{\raggedright\arraybackslash}p{0.30\linewidth}>{\raggedright\arraybackslash}p{0.18\linewidth}>{\raggedright\arraybackslash}p{0.19\linewidth}>{\raggedright\arraybackslash}p{0.20\linewidth}}
\toprule
Statement & Mathematical validity & S6-AE applicability & Scientific novelty \\
\midrule
Per-niche $2\epsilon$ quality & PROVED conditional & PARTIAL & STANDARD \\
One-step $2\epsilon+2\delta$ tracking & PROVED conditional & PARTIAL/OPEN & STANDARD \\
Historical regret decomposition & PROVED conditional & OPEN & STANDARD \\
Historical regret recurrence and sum & PROVED conditional & OPEN & STANDARD \\
Bounded long-run tracking regret & OPEN; not implied & OPEN & NOT CLAIMED \\
Tracking implies reachability/training value & DISPROVED without alignment & NOT APPLICABLE as a guarantee & NOT CLAIMED \\
\bottomrule
\end{tabular}
\end{center}

Approximate maximization under uniform value error is elementary.  Adding a
bounded objective drift and telescoping the resulting regret is in the
standard spirit of dynamic-regret and moving-optimum analyses
\cite{Zinkevich2003,Dang2017,OlivetoZarges2015}.  Fixed-cell elite retention is
also standard QD machinery \cite{MouretClune2015,Pugh2016}.  The only
BLADE-specific ingredients used here are the exact current-HGB niche elite,
parent--child union, and $P=T$, which make the previous elite available for
one recursive step.  A potentially novel result would need a non-vacuous,
implementation-supported error/drift certificate or a hardness--reachability--
training-value relation specific to FOGA-11 and BQ adaptation.  This stage
establishes none of those.

\subsection{Stage 3A research decision}
The proved inequalities are internally valid, but their decisive premises
are not currently supported by S6-AE: exact current LBQ is unavailable on the
complete candidate union, and the historical recurrence requires drift bounds
for discarded candidates across the entire latent ledger.  Counterexamples
block any replacement of those premises by coverage, perfect present-time
ranking alone, or low tracking regret.

\textbf{C -- TRACKING\_ASSUMPTIONS\_NOT\_JUSTIFIED.}
